\documentclass[10pt, a4paper]{article}

\usepackage{geometry}
\geometry{left=2.54cm, right=2.54cm, bottom=2.54cm, top=1.27cm, includefoot}

\usepackage{fontspec}
\setmainfont{Georgia}
\usepackage{setspace}
\singlespacing 
\usepackage{microtype} 

\usepackage{titlesec}
\titleformat{\section}{\normalfont\fontsize{13}{15}\bfseries}{\thesection}{1em}{}
\titleformat{\subsection}{\normalfont\fontsize{11}{13}\bfseries\itshape}{\thesubsection}{1em}{}
\titleformat{\subsubsection}{\normalfont\fontsize{10}{12}\bfseries}{\thesubsubsection}{1em}{}

\usepackage{caption}
\captionsetup{font={bf, normalsize}, labelsep=period, name=Figure, justification=centering}

\usepackage[backend=biber, style=apa]{biblatex}
\addbibresource{references.bib} 

\usepackage{graphicx}
\usepackage{booktabs} 

\begin{document}

\begin{titlepage}
    \centering
    {\fontsize{14}{16}\selectfont Hochschule für angewandte Wissenschaften München\par}
    {\fontsize{14}{16}\selectfont Fakultät für Informatik und Mathematik\par}
    \vspace{3cm}
    {\fontsize{20}{24}\bfseries\selectfont Titel der Arbeit\par}
    \vspace{1.5cm}
    {\fontsize{14}{16}\bfseries\selectfont Bachelorarbeit\par}
    \vspace{2cm}
    \begin{flushleft}
        \normalsize
        \textbf{Eingereicht von:}\\
        $[$Vorname Name$]$\\
        \textbf{Studiengang:} Betriebswirtschaft\\
        \textbf{Matrikelnummer:} xxxxxxxx\\
        \textbf{Kontaktdaten:} $[$alias$]$@hm.edu\\
        \vspace{1cm}
        \textbf{Erstprüfer:} Prof. Dr. Markus Thimmel\\
        \textbf{Zweitprüfer:} Prof. Dr. Christoph Ertl\\
        \textbf{Abgabedatum:} xx.xx.xxxx
    \end{flushleft}
\end{titlepage}

\newpage
\section*{Ggf. Sperrvermerk}
Die vorliegende Arbeit beinhaltet interne vertrauliche Informationen der Firma $[$Alpha GmbH$]$.\par
Die Weitergabe des Inhalts der Arbeit im Gesamten oder in Teilen sowie das Anfertigen von Kopien und Abschriften – auch in digitaler Form – sind grundsätzlich untersagt.\par
Ausnahmen bedürfen der schriftlichen Zustimmung der Firma.

\vspace{2cm}
\section*{Ggf. General Disclaimer KI-Nutzung}
Die vorliegende Arbeit wurde mithilfe Künstlicher Intelligenz auf Grammatik, Stil und Rechtsschreibung überprüft. Der inhaltliche Sinn sowie die wissenschaftlichen Aussagen/Interpretation des Textes wurden dabei nicht verändert.\par
\textbf{Beispiel aus dem Text:}\par
\textit{Vorher:} Die Auswertung von den Interviews zeigt, das mittelständische Unternehmen mit vielen modernen betrieblichen Anforderungen konfrontiert sind, die mehr Einfluss auf die Ausgestaltung und Nutzung der Kostenrechnung haben.\par
\textit{Nachher:} Die Auswertung der Interviews zeigt, dass mittelständische Unternehmen mit einer Reihe moderner betrieblicher Anforderungen konfrontiert sind, die sich zunehmend auf die Ausgestaltung und Nutzung der Kostenrechnung auswirken.

\newpage
\section*{Abstract / Kurzfassung}
$[$Your abstract goes here.$]$

\newpage
\tableofcontents
\newpage

\section{Einleitung}
Bitte halten Sie sich bei der Erstellung Ihrer Seminararbeit an die folgenden grundlegenden Formatierungsrichtlinien. Der gesamte Text ist im Blocksatz zu setzen. \textcite{alridhawi2020} discusses AI challenges, while other works cover IoT networks \parencite{elayan2021}.

\section{Abbildungen und Tabellen}
Die Begriffe Figure und Table sind immer auszuschreiben.

\begin{table}[hbt!]
\centering
\caption{A Very Nice Table}
\label{tab:nice_table}
\begin{tabular}{lcc}
\toprule
 & \textbf{Treatment 1} & \textbf{Treatment 2} \\
\midrule
Setting A & 125 & 95 \\
Setting B & 85 & 102 \\
Setting C & 98 & 85 \\
\bottomrule
\end{tabular}
\end{table}

\newpage
\printbibliography

\end{document}
